\bookchapter{Capítulo 13: Apps de herramientas recomendadas}

\chapterintro{Este capítulo reúne herramientas de productividad, salud y viajes, además de funciones prácticas del sistema y entre dispositivos: limpieza y antispam tipo Qingzhuo, bloqueo integrado del teléfono, clima y alertas sísmicas; control remoto del PC con Sunlogin/ToDesk, oficina con IA tipo Trae, fondos animados, envío rápido en red local con LocalSend, portapapeles entre dispositivos, códigos dinámicos con Aegis/Microsoft Authenticator, agenda con TickTick y escaneo de documentos en papel.}

\booksection{I. Herramientas de productividad: gestiona tu tiempo y tu atención}

Las apps de productividad son el «esqueleto» de tu vida digital: te ayudan a poner orden en ideas, tareas y archivos dispersos. Aquí van recomendaciones por situación.

\booksubsection{Listas de tareas (gestión de pendientes)}

\begin{itemize}
  \item TickTick Público: quien quiera unificar trabajo, vida y estudio en un mismo calendario de tareas; una de las recomendaciones clave de esta sección. Puntos fuertes: sincronización entre móvil/PC/tablet/web; lista + calendario + Pomodoro en una sola app; entrada en lenguaje natural («reunión mañana a las 3», «informe semanal cada lunes») que reconoce fechas y repeticiones; subtareas, etiquetas, prioridades, adjuntos y listas compartidas. Ideal para combinar planificación, seguimiento diario y revisión.
  \item Cómo sacarle más partido a TickTick: Bandeja de entrada: manda todo a «Hoy / Bandeja» y por la noche repártelo entre trabajo, familia y estudio.
  \item Vista de calendario: arrastra reuniones y fechas límite al calendario para ver choques de un vistazo; es más claro que una simple lista.
  \item Tareas repetitivas: medicación, pagos o gimnasio como tareas con reglas, sin tener que rellenar cada vez.
  \item Pomodoro: inicia la concentración desde la propia tarea y mide cuánto has «terminado de verdad» esta semana.
  \item Integración con el teléfono: el asistente de voz con «recuérdame…» crea recordatorios del sistema; para fechas límite importantes, actívalo también en el calendario del sistema como red de seguridad.
\end{itemize}

Todoist Público: usuarios minimalistas y multiplataforma. Puntos fuertes: interfaz sobria y clara, jerarquía de proyectos (tarea subtarea subtarea) y sistema de puntos «karma» para crear hábitos.

Microsoft To Do Público: oficinistas y usuarios de Outlook. Puntos fuertes: gratuito y sin anuncios; la función «Mi día» se vacía cada día y alivia la carga mental; gran integración con el ecosistema Microsoft.

Notas / Recordatorios del sistema Público: usuarios ligeros que no quieren instalar apps de terceros. Puntos fuertes: experiencia nativa, cero retrasos, se añaden con Siri o el asistente de voz, y se integran con los widgets del sistema.

\booksubsection{Temporizador Pomodoro (entrenar la concentración)}

\begin{itemize}
  \item Forest Público: a quien se le va la mano con el móvil. Puntos fuertes: plantas un árbol mientras te concentras; si sales antes de tiempo, el árbol se seca. Con esa «vida virtual» crea refuerzo positivo y puedes plantar con amigos.
  \item Tide Público: quien necesita ruido blanco para concentrarse o dormir. Puntos fuertes: sonidos naturales (lluvia, olas, bosque) + Pomodoro + meditación en una sola app, con interfaz minimalista y relajante.
  \item Principio: 25 minutos de concentración + 5 de descanso; puedes ajustarlo a 50/10 o 90/20 según tu estado.
\end{itemize}

\booksubsection{Calendario (visualizar el tiempo)}

\begin{itemize}
  \item Calendario del sistema Público: todos. Puntos fuertes: notificaciones nativas muy estables y muy bajo consumo; se puede añadir con Siri o el asistente de voz.
  \item Google Calendar Público: usuarios con varios dispositivos o que comparten agenda. Puntos fuertes: la función «Objetivos» te organiza sola tiempo para deporte o estudio, y con Gmail reconoce eventos en los correos.
\end{itemize}

\booksubsection{Notas (consolidar conocimiento)}

\begin{itemize}
  \item Notion Público: quien quiera gestionar notas, tareas, bases de datos y wiki en un solo lugar; estudiantes, freelancers y equipos. Puntos fuertes: base de conocimiento todoterreno, edición por bloques, páginas personales, listas de lectura, tableros de proyectos y un rico ecosistema de plantillas.
  \item Obsidian Público: gestores de conocimiento avanzados, investigadores y escritores. Puntos fuertes: almacenamiento local en Markdown, enlaces bidireccionales y grafo de relaciones para convertir tus notas en un «segundo cerebro»; los datos son totalmente tuyos.
  \item flomo Público: quien quiere capturar ideas sueltas sin escribir textos largos. Puntos fuertes: notas tipo tarjeta, como publicar en redes, con etiquetas \# y línea de tiempo para recuperar lo anterior.
  \item Notas de Apple Público: usuarios de iPhone/Mac con necesidades de registro ligeras. Puntos fuertes: abre rápido, admite escanear documentos, dibujar, listas y sincronización con iCloud; gratis y muy integrado con el sistema.
  \item Evernote Público: usuarios avanzados que recortan muchas páginas web y archivan material. Puntos fuertes: potente extensión de recorte web, OCR de texto en imágenes, muy estable con años de trayectoria; la versión de pago es de suscripción.
\end{itemize}

\booksubsection{Herramientas PDF}

\begin{itemize}
  \item PDF Expert Público: usuarios de iPad/iPhone, profesionales que anotan contratos y estudiantes. Puntos fuertes: considerado el mejor editor de PDF en iOS; anotaciones fluidas, firmas, combinar/dividir y rellenar formularios.
  \item WPS Público: oficinistas que manejan Word/Excel/PPT/PDF a la vez. Puntos fuertes: suite de oficina todo en uno, buena compatibilidad con nubes chinas y archivos de WeChat; la versión gratuita cubre el uso diario.
  \item Adobe Acrobat Público: quien necesita un formato PDF muy normativo (legal, diseño, edición). Puntos fuertes: estándar oficial, conversión fiable a PDF/A, sincronización entre dispositivos y firma electrónica profesional.
\end{itemize}

\booksubsection{Escaneo (digitalizar documentos en papel)}

Contratos, facturas, deberes, pizarras, copias del DNI… la cámara del móvil ya sustituye a la mayoría de escáneres. Flujo básico: foto nítida recorte y enderezado automático exportar a PDF/imagen OCR opcional.

\begin{itemize}
  \item CamScanner Público: profesionales y estudiantes que escanean contratos, recibos o documentos con frecuencia. Puntos fuertes: detección automática de bordes, corrección de perspectiva, OCR y creación de PDF multipágina con copia en la nube.
  \item Microsoft Lens Público: estudiantes y quien quiera convertir pizarras o páginas de libro en Word. Puntos fuertes: gratis y sin anuncios, exporta a OneNote/Word con un clic y reconoce texto con mucha precisión.
  \item Escaneo en Notas de Apple Público: usuarios de iPhone con necesidades de escaneo ligeras. Puntos fuertes: no hace falta instalar nada; con el icono de cámara de Notas escaneas, se quitan sombras y puedes firmar; gratis y muy cómodo.
  \item Baimiao / CamScanner web / «Escaneo de documentos» de la cámara de marca: Huawei, Xiaomi, OPPO y vivo suelen incluir un modo «documento/escaneo» en la cámara de fábrica; en Android también sirve Google Camera o el escáner del sistema. Ideal para una foto urgente de una página de contrato.
  \item Trucos de foto: luz uniforme, móvil paralelo al papel y encuadre en las cuatro esquinas; en documentos de varias páginas, fotografía en orden y luego combina en un PDF; con documentos sensibles, cuida las marcas de agua y dónde los guardas; cuando termines, bórralos también de «Eliminados recientemente».
\end{itemize}

\booksection{II. Salud y ejercicio}

Las apps de salud son el «asistente personal» más infravalorado del móvil: siguen tu sueño, ritmo cardíaco y actividad de forma continua, y reflejan tu estado a largo plazo mejor que una visita esporádica al médico.

\booksubsection{Deporte y fitness}

\begin{itemize}
  \item Keep Público: entrena en casa, principiantes o quien quiera perder grasa y tonificar. Puntos fuertes: muchísimos cursos gratuitos, de HIIT a yoga, planes de entrenamiento generados por IA y una comunidad muy activa.
  \item Nike Training Club (NTC) Público: quien busca guía de entrenadores profesionales y un plan sistemático. Puntos fuertes: app oficial de Nike, todos los cursos gratuitos, demostraciones grabadas por entrenadores y clasificación clara por grupos musculares y material.
  \item Apple Fitness Público: usuarios de Apple Watch que quieren motivarse con los anillos de actividad. Puntos fuertes: tres anillos diarios (moverse/ejercicio/de pie), objetivos visuales y datos de ritmo cardíaco del reloj sin fisuras.
\end{itemize}

\booksubsection{Seguimiento del sueño}

\begin{itemize}
  \item AutoSleep Público: usuarios de Apple Watch que quieren saber su proporción de sueño profundo/ligero. Puntos fuertes: hace seguimiento automático al dormir con el reloj, sin iniciar ni detener a mano, y ofrece una puntuación de calidad y consejos para acostarse.
  \item Sleep Cycle Público: quien no duerme con reloj y quiere monitorizar con el móvil. Puntos fuertes: analiza con el micrófono los movimientos y los ronquidos, y te despierta en fase de sueño ligero para un levantarte menos duro.
  \item Mi Fitness Público: usuarios de pulseras o relojes Xiaomi. Puntos fuertes: gratis y muy integrado con la pulsera; seguimiento de ritmo cardíaco, oxígeno y sueño en una sola app, con excelente relación calidad-precio.
\end{itemize}

\booksubsection{Meditación y relajación}

\begin{itemize}
  \item Calm Público: oficinistas con estrés o a quien le cuesta dormir. Puntos fuertes: meditación guiada «Daily Calm» + historias para dormir narradas por famosos; interfaz muy cuidada y mucho contenido en inglés.
  \item Tide Público: usuarios de habla china que necesitan ruido blanco para concentrarse o dormir. Puntos fuertes: producto chino de calidad, con sonidos naturales, cursos de meditación y Pomodoro; las guías en chino resultan más cercanas.
  \item Little Sleep Público: quien necesita sonidos para dormir o grabar sonámbulos/ronquidos. Puntos fuertes: enorme variedad de ruidos blancos combinables (lluvia + ventilador + fogata), apagado programado al dormir y grabación de sonámbulos, útil y divertida.
\end{itemize}

\booksubsection{Salud femenina}

\begin{itemize}
  \item Flo Público: quien registra el ciclo, busca embarazo o quiere cuidar la salud femenina. Puntos fuertes: predicción con IA de la próxima regla y la ovulación, con buena precisión; comunidad con preguntas anónimas y diseño muy pulido.
  \item Meiyou Público: usuarias en China que prefieren comunidad y divulgación en chino. Puntos fuertes: contenido local muy completo, con conocimientos por etapas: ciclo, búsqueda de embarazo, gestación y crianza.
\end{itemize}

\booksubsection{Recordatorios de medicación}

\begin{itemize}
  \item Asistente de medicación. Para quién: personas con enfermedades crónicas que toman medicación a largo plazo e hijos que cuidan la medicación de sus padres mayores. Punto fuerte: identifica medicamentos escaneando el código de barras, crea automáticamente la lista de medicación, avisa con ventanas emergentes a la hora y registra las tomas para compartirlas con la familia.
\end{itemize}

\booksubsection{Centro de datos de salud}

\begin{itemize}
  \item Salud de Apple / Google Fit Público: todos los usuarios de smartphones. Puntos fuertes: reúne los datos del móvil, la pulsera y las apps de salud de terceros (pasos, ritmo cardíaco, sueño, ciclo, etc.) en un único «expediente de salud».
\end{itemize}

\booksection{III. Viajes y desplazamientos}

Fuera de casa, unas pocas apps bien elegidas te ahorran un montón de problemas. Aquí van recomendadas por etapa del viaje.

\booksubsection{Vuelos y trenes}

\begin{itemize}
  \item Umetrip. Para quién: profesionales que vuelan con frecuencia. Punto fuerte: datos oficiales de aviación civil; es la fuente más fiable para el estado de los vuelos, los cambios de puerta y los avisos de retraso, y permite consultar el modelo de avión y el mapa de asientos.
  \item 12306 Público: todos los que viajan en tren en China. Puntos fuertes: el único canal oficial de compra de confianza; admite lista de espera (compra automática si hay devoluciones) y trámite sencillo de billetes de estudiante o de grupo.
  \item VariFlight. Para quién: quienes son sensibles a la puntualidad de los vuelos o necesitan llevar o recoger a alguien del aeropuerto. Punto fuerte: predice con precisión los retrasos y permite compartir el itinerario con la familia para que vea en tiempo real tu ubicación y tu hora de llegada.
\end{itemize}

\booksubsection{Mapas y navegación}

\begin{itemize}
  \item Amap (Gaode) Público: usuarios urbanos en China, conductores y quienes piden taxi. Puntos fuertes: tráfico en tiempo real muy preciso, buen desvío de atascos, taxi y carpooling en la misma app y muchas opciones de transporte público.
  \item Baidu Maps Público: quien necesita información de lugares y navegación interior. Puntos fuertes: mapa panorámico y navegación AR, planos detallados de centros comerciales y atracciones, y muchos paquetes de voz.
  \item Google Maps Público: viajeros al extranjero. Puntos fuertes: la cobertura más completa del mundo, reseñas locales de fiar; Street View y mapas sin conexión son un salvavidas lejos de casa.
\end{itemize}

\booksubsection{Traducción}

\begin{itemize}
  \item Google Traductor Público: viajeros y estudiantes de idiomas. Puntos fuertes: traducción conversacional bidireccional en tiempo real y traducción instantánea con la cámara; imprescindible de viaje, con descarga de idiomas sin conexión.
\end{itemize}

\booksubsection{Conversión de divisas}

\begin{itemize}
  \item Jijian Huilü (conversor minimalista de divisas). Para quién: quienes viajan al extranjero o compran online en tiendas internacionales. Punto fuerte: interfaz muy limpia, más de 160 monedas, tipos de cambio actualizados en tiempo real, comparación de varias monedas a la vez y conversión rápida manteniendo pulsado el número.
\end{itemize}

\booksubsection{Clima}

\begin{itemize}
  \item La app del tiempo del sistema / China Weather (weather.com.cn). Para quién: todo el mundo; es la primera opción para consultar el pronóstico y las alertas meteorológicas oficiales. Punto fuerte: las apps del tiempo de los fabricantes están conectadas con canales oficiales como «China Weather», así que las alertas de lluvias torrenciales, tifones, olas de calor, etc. llegan como banners o notificaciones, con gran autoridad.
  \item Caiyun Weather Público: quien necesita previsión de lluvia minuto a minuto. Puntos fuertes: precisión líder en «¿va a llover en las próximas dos horas?», mapa dinámico de lluvia muy visual y avisos proactivos antes de que empiece a llover.
  \item Carrot Weather Público: entusiastas del clima que buscan diversión y widgets ricos. Puntos fuertes: locutor de IA con humor ácido, muchos widgets personalizables y radar profesional; la versión de pago ofrece más funciones.
  \item Clave de alertas: las apps de clima necesitan permiso de notificaciones para recibir avisos; las alertas tienen cuatro niveles —azul, amarillo, naranja y rojo— y el rojo indica el mayor riesgo de fenómenos peligrosos, así que ajusta tus planes según el aviso. Más formas de activarlas en la sección 7 de este capítulo «Alertas meteorológicas / desastres naturales / sismos».
\end{itemize}

\booksubsection{Guías de viaje}

\begin{itemize}
  \item Mafengwo Público: a quien le gusta planear en detalle, viajes por carretera o turismo de profundidad. Puntos fuertes: relatos originales de usuarios muy completos, guías de destino, rutas y consejos para evitar problemas, y descarga de guías sin conexión.
  \item Xiaohongshu (RED) Público: jóvenes viajeros que buscan experiencias reales, rincones poco conocidos y buena comida. Puntos fuertes: formato visual de imagen y vídeo corto, contenidos como «3 días en XX» muy actualizados y buen motor de búsqueda.
\end{itemize}

\booksection{IV. Introducción a la automatización}

La automatización es una palanca: un poco de configuración a cambio de ahorrar esfuerzo durante mucho tiempo. Puede convertir en cero las decenas de acciones repetidas de cada día. Aquí solo damos el concepto; el detalle y los tutoriales están en el Capítulo 23: Pensamiento computacional e introducción a la automatización.

\begin{itemize}
  \item Atajos de iOS: la herramienta oficial de automatización de Apple, sin programar. En un solo gesto puedes encadenar acciones como «al llegar a casa: activar WiFi + bajar el brillo + anunciar el clima».
  \item Automatización en Android: las escenas del sistema (por ejemplo «Tareas automáticas» de Xiaomi o «Inteligencia de escena» de Huawei) son las más fáciles; Tasker es muy potente pero tiene una curva de aprendizaje pronunciada, ideal para usuarios avanzados.
  \item Idea clave: primero detecta «lo que repites cada día» y luego piensa «¿puedo sustituirlo por un disparador + varias acciones?». Ejemplos detallados y usos avanzados en el Capítulo 23.
\end{itemize}

\booksection{V. Navegadores y buscadores: tu puerta de entrada a la información}

El navegador es tu ventana a internet y el buscador, la brújula para orientarte en el océano de información. Usarlos bien te evita muchos rodeos.

\booksubsection{Cómo funciona un navegador:}

Cuando escribes una dirección, el navegador: (1) consulta al servidor DNS la IP del dominio; (2) envía una petición HTTP/HTTPS a esa IP; (3) el servidor devuelve el código de la página (HTML, CSS, JavaScript); (4) el navegador interpreta el código y dibuja la página que ves.

\booksubsection{Navegadores recomendados:}

\begin{itemize}
  \item Microsoft Edge: basado en Chromium, rápido y muy compatible. Trae el asistente Copilot con IA para resumir páginas, traducir y conversar. Sincroniza entre dispositivos; recomendación personal como navegador principal.
  \item Safari: navegador por defecto en dispositivos Apple, muy integrado con el sistema y eficiente en batería. Buena privacidad, aunque con menos extensiones.
  \item Chrome: el más usado del mundo y con más extensiones. Consume bastante memoria y en el móvil gasta más batería.
\end{itemize}

\booksubsection{Qué buscador elegir:}

\begin{itemize}
  \item Bing: buscador de Microsoft, muy integrado con Edge. Sus resúmenes con IA dan la respuesta directa, no solo enlaces, lo que agiliza mucho la búsqueda. La búsqueda de imágenes también está bien.
  \item Yandex: el mayor buscador de Rusia, muy fuerte en búsqueda de imágenes y búsqueda inversa. Muy útil para encontrar fotos similares o identificar objetos.
  \item Baidu: el principal buscador en China; tiene bastantes anuncios, así que conviene criterio para filtrar.
  \item Google: el buscador más preciso del mundo, pero en China requiere una red especial para acceder.
\end{itemize}

\booksubsection{Guía de búsqueda eficiente:}

\begin{itemize}
  \item Aprovecha la búsqueda con IA: los resúmenes de Bing y el Copilot de Edge responden preguntas y destilan la información clave entre miles de resultados; mucho más eficaz que ir abriendo enlaces uno a uno.
  \item Detectar desinformación y anuncios: las entradas marcadas como «Anuncio» o «Patrocinado» son publicidad de pago, no resultados orgánicos. Fíjate en la autoridad de la fuente y prioriza sitios oficiales y medios reconocidos.
  \item Descargar apps de forma correcta: escribe en el buscador «nombre de la app + sitio oficial», por ejemplo «WeChat sitio oficial», y entra en resultados con dominio oficial. No instales desde tiendas de terceros: suelen incluir software no deseado.
  \item Cómo reconocer el sitio oficial: para Windows 11 debes ir a microsoft.com; para Chrome, a google.com. Comprobar que el dominio coincide con la marca es el método más básico.
\end{itemize}

\booksubsection{Conocimientos sobre dominios web:}

El dominio es la «dirección» del sitio; saber leerlo te ayuda a juzgar su naturaleza y origen:

\begin{itemize}
  \item Significado de los sufijos más comunes: .com — empresa comercial, el más habitual
  \item .cn — dominio nacional de China
  \item .org — organización sin ánimo de lucro
  \item .net — servicios de red
  \item .edu — instituciones educativas
  \item .gov — organismos gubernamentales
  \item .xyz/.tech/.info — dominios genéricos nuevos, de acceso fácil; precaución
\end{itemize}

Dominios por país/región: .jp — Japón

.kr — Corea del Sur

.uk — Reino Unido

.de — Alemania

.tw — Taiwán

.hk — Hong Kong

Aviso de seguridad: los sitios .gov y .edu suelen ser más fiables porque el registro es más exigente. Dominios baratos como .xyz o .top se usan mucho en sitios de phishing; al visitarlos, extrema la precaución.

\booksubsection{Bloqueo de anuncios:}

\begin{itemize}
  \item AdGuard Público: a quien le molesten ventanas emergentes, anuncios superpuestos o publicidad dentro de apps. Puntos fuertes: bloqueo a nivel de sistema (no solo en el navegador), filtra banners y rastreo de pantallas de inicio, ahorra datos y acarga las páginas; la versión de pago cubre todos tus dispositivos.
\end{itemize}

\booksubsection{Navegadores centrados en la privacidad:}

\begin{itemize}
  \item Brave Público: navegadores del día a día que valoran la privacidad y no quieren ser rastreados. Puntos fuertes: basado en Chromium, bloquea anuncios y rastreadores desde el primer momento, trae modo de navegación con Tor, va más rápido que Chrome y usa menos memoria.
  \item DuckDuckGo Público: quien no quiere que al buscar una palabra le persigan anuncios por toda la red. Puntos fuertes: incluye su propio buscador (sin guardar el historial), botón «Fire» para borrar todo rastro de un clic y experiencia uniforme entre plataformas.
\end{itemize}

\booksubsection{Comparativa de navegadores:}

\begin{center}
\begin{tabularx}{\linewidth}{|X|X|X|}
\hline
\textbf{Navegador} & \textbf{Punto fuerte} & \textbf{Público} \tabularnewline
\hline
Microsoft Edge & Núcleo Chromium + Copilot con IA: resume páginas, traduce y conversa & Oficinistas orientados a la eficiencia y usuarios intensivos de IA (recomendado como principal) \tabularnewline
Safari & Muy integrado con el ecosistema Apple, eficiente en batería y con buena privacidad & Usuarios de iPhone/Mac que cuidan la autonomía y la integración \tabularnewline
Chrome & El más usado del mundo y con el ecosistema de extensiones más rico & Usuarios avanzados que necesitan muchas extensiones o sincronización multi-dispositivo \tabularnewline
Brave & Bloquea anuncios y rastreadores desde el primer momento y va más rápido & Quien prioriza la privacidad y quiere bloquear anuncios sin instalar extensiones \tabularnewline
DuckDuckGo & Su buscador no guarda historial y se limpia todo con un clic & Usuarios sensibles a la privacidad que no quieren seguimiento personalizado \tabularnewline
\hline
\end{tabularx}
\end{center}

\begin{tcbpractice}{Crea una automatización «para relajarte antes de dormir»}
\begin{enumerate}
  \item Define el disparador: una hora concreta (por ejemplo, 22:30).
  \item Añade las acciones: bajar el volumen > activar No molestar > reproducir la lista de relajación > bajar el brillo.
  \item Guarda y prueba.
\end{enumerate}

\end{tcbpractice}

\booksection{VI. Herramientas prácticas del sistema: limpieza, antispam y seguridad}

Más allá de las apps de productividad o salud, lo que marca el día a día suele ser: «que no te molesten, que haya espacio suficiente y que los estafadores no entren». Aquí repasamos software útil y luego detallamos el bloqueo de spam del propio teléfono: muchos llevan años con un móvil sin haber activado nunca estos interruptores.

\booksubsection{Limpieza del sistema y optimización del almacenamiento}

\begin{itemize}
  \item Qingzhuo Público: usuarios de Android que quieren limpiar caché, archivos residuales y carpetas vacías sin el ruido de los «gestores de teléfono» comerciales. Puntos fuertes: ocupación pequeña, sin anuncios, análisis profundo de residuos, localización de archivos grandes y diagnóstico de uso de almacenamiento; más discreto que los «acelerar de un toque» comerciales, ideal para mantenimiento periódico si tienes mano. AtenCIÓN: antes de limpiar, revisa las rutas; no selecciones a ciegas WeChat, galería o documentos; para caché del sistema, prioriza la limpieza de la marca y usa herramientas de terceros solo como complemento.
  \item Limpieza de almacenamiento del sistema Público: todos, sobre todo personas mayores o quienes no quieren complicaciones. Puntos fuertes: Android [Ajustes] > [Almacenamiento] > [Limpiar y acelerar]; iOS [Ajustes] > [General] > [Almacenamiento del iPhone]. El sistema conoce mejor la caché de sus apps y tiene el menor riesgo de borrar algo importante: debe ser la primera opción.
  \item Gestión de almacenamiento en WeChat/QQ Público: quienes chatean mucho. Puntos fuertes: WeChat [Yo] > [Ajustes] > [General] > [Almacenamiento]; permite limpiar por conversación o tipo de archivo, y es más seguro que las herramientas externas de «limpiar WeChat».
  \item Principio: si el sistema o la app ya lo trae, no uses terceros; si puedes elegir por ruta o tipo, no hagas un «limpiar todo». El objetivo es liberar espacio, no llevarte por delante datos importantes.
\end{itemize}

\booksubsection{Apps de antispam y seguridad}

\begin{itemize}
  \item Centro Nacional Anti-Fraude Público: todos, en especial personas mayores, personal financiero y quienes compran en línea o atienden llamadas de desconocidos. Puntos fuertes: app oficial con alertas de llamadas, avisos de riesgo en SMS y detección de URLs sospechosas; conectada con las bases antifraude de la policía, es el eslabón con más respaldo oficial de las «tres líneas de defensa».
  \item Centro de seguridad de marca / gestor del teléfono Público: usuarios de esa marca. Puntos fuertes: Huawei Phone Manager, Xiaomi Security Center, OPPO Phone Manager, vivo Phone Manager, etc., con bloqueo de spam, antivirus, permisos y protección de pagos; suelen ser gratuitos y mejor integrados con el sistema. No instales un segundo paquete comercial que pueda chocar en permisos.
  \item Servicio antispam del operador Público: usuarios bombardeados por llamadas automáticas o voz con IA. Puntos fuertes: China Mobile, envía KTFSR al 10086; China Telecom, sigue «Tianyi Antispam»; China Unicom, configúralo en «Unicom Mobile Manager». El bloqueo ocurre en la red del operador y corta un buen número de llamadas sin que suene el teléfono.
  \item AdGuard Público: a quien le molesten ventanas emergentes, anuncios superpuestos o publicidad dentro de apps. Puntos fuertes: bloqueo a nivel de sistema (no solo en el navegador), filtra banners y rastreo de pantallas de inicio, ahorra datos y acarga las páginas; la versión de pago cubre todos tus dispositivos.
  \item Cuidado con los paquetes «limpiar / acelerar / bloquear spam»: algunos gestores de terceros se quedan en segundo plano y piden permisos de contactos y SMS una y otra vez; su bloqueo no siempre supera al del sistema y suman consumo y riesgos de privacidad. Prioriza las funciones del sistema y, si acaso, la app oficial antifraude.
\end{itemize}

\textbf{Principios para elegir software útil}\par

\begin{itemize}
  \item Primero el sistema: para bloquear, limpiar o gestionar permisos, activa antes los interruptores del propio teléfono.
  \item Complemento oficial: el Centro Anti-Fraude Nacional y el antispam del operador son canales oficiales y se pueden sumar.
  \item Terceros con mesura: basta una app por función para evitar choques de permisos y peleas por recursos en segundo plano.
  \item Menos es más: cuantas más herramientas instalas, más notificaciones y permisos, y más probable que te molesten las propias herramientas.
\end{itemize}

\booksubsection{Bloqueo de spam del propio teléfono: cómo activarlo según la marca}

Rutas por marca. Los nombres pueden variar según la versión del sistema; busca en Ajustes «spam», «bloqueo» o «llamadas desconocidas». Complementa el Capítulo 8 «Gestión de llamadas spam»: aquí nos centramos en dónde están los interruptores; allí, en el etiquetado de números y reclamos por errores.

\begin{center}
\begin{tabularx}{\linewidth}{|X|X|X|}
\hline
\textbf{Plataforma} & \textbf{Ruta de acceso} & \textbf{Opciones recomendadas} \tabularnewline
\hline
iPhone / iOS & {}[Ajustes] > [Teléfono] / [Mensajes] & Silenciar llamadas desconocidas; contactos bloqueados; «Filtrar remitentes desconocidos» en Mensajes \tabularnewline
Xiaomi / Redmi & App Teléfono > Ajustes > Bloqueo de spam; o Centro de seguridad & Bloqueo inteligente, colgar llamadas de un solo tono de desconocidos, filtro de palabras clave en SMS \tabularnewline
Huawei / Honor & Teléfono > Bloqueo de spam; Gestor del teléfono > Bloqueo de spam & Identificación de llamadas, lista negra, bloqueo inteligente de SMS \tabularnewline
OPPO / realme / OnePlus & App Teléfono > Bloqueo de spam; o Este dispositivo > Bloqueo de spam & Bloqueo de llamadas, identificación de estafas, reglas de bloqueo de SMS \tabularnewline
vivo / iQOO & Teléfono > Bloqueo de llamadas spam; i Manager & Bloqueo de spam, alertas de estafas, filtrado de mensajes \tabularnewline
Samsung & Teléfono > Ajustes > Bloquear números / Identificador de llamadas y spam & Bloquear números desconocidos, identificación de spam (según la región) \tabularnewline
\hline
\end{tabularx}
\end{center}

\begin{tcbinfo}{Cómo combinar las tres capas de protección}
Recomendación diaria: «bloqueo del sistema + antispam del operador + Centro Anti-Fraude Nacional». El sistema se encarga del etiquetado y la lista negra; el operador corta llamadas automáticas en la red; el centro antifraude aporta alertas oficiales. No hace falta llevarlo todo al máximo, o bloquearás también a repartidores, comida o médicos: deja pasar tus contactos y los recientes.

\end{tcbinfo}

\begin{tcbpractice}{Práctica: activa hoy el bloqueo de spam}
\begin{enumerate}
  \item Abre los ajustes de la app Teléfono, busca «Bloqueo de spam / Silenciar llamadas desconocidas» y activa el bloqueo inteligente.
  \item Revisa también el móvil de tus padres: misma ruta + antispam del operador (en China Mobile, envía KTFSR al 10086).
  \item Instala «Centro Nacional Anti-Fraude» desde la tienda, regístrate y activa las alertas de llamadas.
  \item Prueba con un número de prueba ya etiquetado (o revisa el registro de bloqueos) para confirmar que funciona.
\end{enumerate}

\end{tcbpractice}

\booksection{VII. Alertas meteorológicas, de desastres naturales y sísmicas}

El bloqueo de spam frena «estafas por teléfono»; las alertas, «desastres y clima extremo». Muchos solo instalan una app de clima para ver la temperatura, sin notificaciones ni alerta sísmica. En los segundos o decenas de segundos previos a un temporal, un tifón o un terremoto, el móvil puede regalarte un tiempo precioso para protegerte.

\booksubsection{Alertas meteorológicas: lluvias, tifones y olas de calor}

\begin{itemize}
  \item Prioriza canales oficiales: las alertas de la Administración Meteorológica de China, el Observatorio Central y las oficinas provinciales son las más fiables. Sigue la cuenta oficial «Clima de China» en WeChat o Weibo, o descarga apps oficiales como «Clima de China».
  \item App de clima del sistema: Huawei Weather, Xiaomi Weather, OPPO Weather, vivo Weather, Clima de Apple, etc.; la mayoría envía alertas del observatorio local. Entra en los ajustes de la app y confirma que «Alertas meteorológicas / clima severo» estén activadas.
  \item Precisión minuto a minuto tipo Caiyun Weather: útil para decidir antes de salir o ir a por los niños si «en esta media hora hace falta paraguas»; las alertas oficiales siguen mandando.
  \item Significado de los colores de alerta (general): Azul: precaución; sal con normalidad y sigue las actualizaciones.
  \item Amarillo: refuerza la prevención y reduce actividades al aire libre innecesarias.
  \item Naranja: severo; quédate en interiores y aléjate de zonas de riesgo.
  \item Rojo: muy severo; sigue las indicaciones de suspensión de clases o trabajo y de evacuación, y no cruces zonas inundadas, torrentes o tramos con viento fuerte.
\end{itemize}

Tipos de alerta habituales: lluvias torrenciales, nevadas intensas, tifones, olas de calor, olas de frío, vientos fuertes, tormentas eléctricas, granizo, hielo en carreteras, niebla densa, sequía y riesgo de deslizamientos.

\booksubsection{Alertas de desastres naturales: no es solo «lluvia»}

\begin{itemize}
  \item Quién las emite: los observatorios, las alertas meteorológicas; los deslizamientos o avenidas suelen publicarse de forma conjunta por recursos naturales, obras hidráulicas y protección civil; los gobiernos locales también avisan por radio de emergencia, SMS o el canal 12379.
  \item Cómo recibirlas por varios canales: Notificaciones del sistema + alertas de la app de clima.
  \item Sigue en WeChat las cuentas oficiales de la «Oficina de Gestión de Emergencias» y la «Agencia Meteorológica» locales.
  \item Canales del sistema nacional de alertas 12379 (SMS o servicio, según la región).
  \item Televisión, radio de emergencia y altavoces de barrio o aldea: en zonas de montaña o en temporada de lluvias no te fíes solo del móvil.
\end{itemize}

Al aire libre y en desplazamientos: con tifón, no te acerques a rótulos ni pases bajo árboles; con lluvias torrenciales, no entres en garajes, túneles o pasos inferiores; en la montaña, ante un aviso de «riesgo meteorológico de deslizamientos», abandona de inmediato laderas y barrancos.

Enlace con el manual de primeros auxilios: si tras el desastre no hay red, se agota la batería o necesitas pedir ayuda con ubicación, consulta el Capítulo 32 Terremotos/desastres y el Capítulo 7 Comunicación por satélite y sin cobertura.

\booksubsection{Alerta sísmica: no es lo mismo que «predecir terremotos»}

La alerta temprana de terremotos (EEW) envía un aviso con segundos o decenas de segundos de antelación a las zonas posiblemente afectadas, aprovechando que las ondas electromagnéticas viajan más rápido que las sísmicas y que las estaciones captan antes las ondas P que las más destructivas ondas S. No predice si habrá terremoto, pero gana un margen para «refugiarse, cortar el fuego o salir del ascensor». Cuanto más cerca del epicentro, menos tiempo hay; a veces no llega ni eso.

\begin{itemize}
  \item Soluciones oficiales y del sector: en China continental son habituales las capacidades de la Red de Alerta de Terremotos de China y servicios impulsados por institutos como el de reducción de desastres de Chengdu; algunos fabricantes preinstalan o permiten recibir alertas en línea.
  \item Apps de terceros: busca en la tienda aplicaciones oficiales o de instituciones de «alerta sísmica» (antes de instalar, revisa el desarrollador y evita copias falsas que piden permisos de más).
  \item Qué hacer al recibir la alerta: en interiores, guarecerte cerca (bajo la mesa, en esquina de muro de carga, junto a muebles sólidos) y protege cabeza y cuello; aléjate de ventanas, fachadas y muebles altos; no corras a ciegas por escaleras ni saltes; si estás en el ascensor, sal en la planta más cercana; en exteriores, busca un lugar abierto y aléjate de cables, rótulos y edificios dañados. Tras el temblor, revisa gas y electricidad y atención a las réplicas.
  \item No silencies ni apagues: permite que la alerta sísmica «traspase» el silencio y el modo No molestar; de noche, deja el móvil al alcance y el volumen del tono sin ponerlo a 0.
\end{itemize}

\booksubsection{Alerta sísmica del propio teléfono: cómo activarla según la marca}

Las rutas pueden variar según la versión del sistema; busca en Ajustes «terremoto», «alerta» o «alerta de seguridad». En iOS el soporte de alertas en China es limitado; suele hacer falta una app de terceros de confianza, con permisos de ubicación y notificaciones activados.

\begin{center}
\begin{tabularx}{\linewidth}{|X|X|X|}
\hline
\textbf{Plataforma} & \textbf{Ruta habitual} & \textbf{Notas} \tabularnewline
\hline
Xiaomi / Redmi & Gestor del teléfono / Centro de seguridad > Alerta sísmica; o busca «terremoto» en Ajustes & Algunas versiones incluyen el servicio; requiere conexión y ubicación \tabularnewline
Huawei / Honor & Ajustes > Seguridad > Notificaciones de emergencia / Alerta sísmica & Puede llamarse «Notificaciones de emergencia» e incluir sismos \tabularnewline
OPPO / realme / OnePlus & Ajustes > Seguridad (SOS / contactos de emergencia) > Alerta sísmica & Según el modelo y la cobertura del servicio en la región \tabularnewline
vivo / iQOO & App de clima o Ajustes > Seguridad/Más > Alerta sísmica & Confirma el menú en tu versión del sistema \tabularnewline
iPhone / iOS & Suele no haber un interruptor completo integrado; instala una app de alerta sísmica de confianza y permite notificaciones & La app necesita ejecutarse en segundo plano y permiso de ubicación; el modo No molestar no debe silenciarla \tabularnewline
Samsung y otros & Ajustes > Seguridad y emergencias / Alertas inalámbricas de emergencia & La cobertura de las «alertas de emergencia» varía por región y operador \tabularnewline
\hline
\end{tabularx}
\end{center}

\begin{tcbwarn}{Errores frecuentes sobre la alerta sísmica}
\begin{itemize}
  \item Alerta $\neq$ predicción: no te dice «si mañana habrá terremoto»; solo avisa cuando ya ha ocurrido y las ondas aún no han llegado.
  \item No en todas partes hay muchos segundos: cerca del epicentro puede haber solo unos instantes, o nada; la alerta beneficia sobre todo a las zonas circundantes.
  \item Apagar las notificaciones es como no tener la app: por buena que sea, si está silenciada, se mata en segundo plano o sin ubicación, puede no llegar.
  \item Apps falsas: desconfía de «predicción de terremotos» que pida contactos o SMS; prioriza lo del sistema o soluciones de instituciones serias.
\end{itemize}

\end{tcbwarn}

\begin{tcbinfo}{Cómo combinar la capacidad de alerta}
Rutina recomendada: «alertas del sistema de clima + canales oficiales de meteorología/emergencias + alerta sísmica (sistema o app de confianza)». Como con el antispam, no hace falta instalar diez apps iguales: lo clave es tener permiso de notificaciones, ubicación permitida y que el modo No molestar no la bloquee. Tras configurar el bloqueo en el móvil de tus padres, revisa también las alertas de clima y sismos.

\end{tcbinfo}

\begin{tcbpractice}{Práctica: revisa tus alertas en 10 minutos}
\begin{enumerate}
  \item Abre la app de clima del sistema, entra en Ajustes y activa «Alertas meteorológicas / clima severo».
  \item Busca en Ajustes del sistema «terremoto» o «alerta» y activa la alerta sísmica/de emergencia (Android); en iPhone, instala una app de alerta sísmica de confianza y permite notificaciones.
  \item Sigue la cuenta oficial de la agencia meteorológica o de emergencias local y permite sus avisos.
  \item Comprueba que el sonido de la alerta no quede anulado por el modo No molestar; repite la revisión en el móvil de tus padres.
\end{enumerate}

\end{tcbpractice}

\booksection{VIII. Colaboración entre dispositivos: control remoto, envío rápido y seguridad de cuentas}

Un móvil fluye mejor cuando «forma equipo» con el ordenador y la tablet. Aquí van herramientas muy usadas: controlar el PC desde el móvil, oficina con IA y trabajo remoto, fondos animados, envío en red local, portapapeles multi-dispositivo, códigos dinámicos 2FA y el enlace entre escaneo en papel y agenda de tareas. La visión más completa de la multi-dispositivo está en el Capítulo 17.

\booksubsection{Controlar el PC desde el móvil: Sunlogin / ToDesk}

\begin{itemize}
  \item Sunlogin Remote Control Público: quien necesita manejar el Windows o Mac de casa o la oficina desde fuera; también se usa mucho para ayudar a personas mayores con el móvil en remoto (Capítulo 27). Puntos fuertes: ecosistema chino maduro; desde la app controlas el escritorio, pasas archivos y enciendes/apagas (según hardware); la versión gratuita suele bastar para uso personal ligero.
  \item ToDesk Público: quien busca estabilidad y baja latencia para teletrabajo o soporte. Puntos fuertes: conexión por código de dispositivo + contraseña, móvil a PC y PC a móvil (según versión), transferencia de archivos y cambio de pantallas cómodo; en China suele ir bien.
  \item Uso básico: instala el cliente en el PC y mantén la sesión instala la app en el móvil inicia sesión con el ID del dispositivo y la contraseña (o la misma cuenta) maneja el escritorio con el trackpad o gestos del móvil.
  \item Límites de seguridad (muy importante): Si un desconocido se hace pasar por «soporte, policía o el operador» y te pide instalar Sunlogin o ToDesk para controlar tu móvil u ordenador, rechaza siempre: es una estafa típica (detalles en el Capítulo 27).
  \item Para uso propio: activa la verificación en dos pasos, cambia la contraseña por defecto y cierra la sesión o desactiva el acceso sin supervisión cuando no lo uses; en el ordenador de la empresa, respeta la política de seguridad.
  \item En lugares públicos, fíjate en que nadie mire la pantalla ni las contraseñas mientras trabajas en remoto.
\end{itemize}

\booksubsection{Oficina con IA: Trae / WorkBuddy / Dumate}

\begin{itemize}
  \item Trae Público: programación y edición de documentos, en entornos chinos. Puntos fuertes: IDE / asistente de programación con IA de ByteDance; el modo SOLO genera un proyecto con una frase; muy útil para código y documentos. Si corre en el PC, desde el móvil puedes seguir tareas con control remoto o las funciones oficiales del producto (según la versión actual). Ver Capítulo 41 Panorama de herramientas de programación con IA.
  \item WorkBuddy Público: profesionales que quieren que la IA ayude con actas, correos, informes y documentos de oficina. Puntos fuertes: asistente de IA orientado a oficina, ideal para tareas de «redactar—pulir—resumir»; se puede iniciar en el móvil y dejar el trabajo pesado al PC.
  \item Dumate Público: equipos o personas que necesitan colaboración y gestión de conocimiento en una sola herramienta. Puntos fuertes: oficina colaborativa con IA, énfasis en reunir material y edición conjunta; en el móvil sirve para ver y editar a medias; el maquetado complejo, mejor en pantalla grande.
  \item Qué significa «controlar el PC desde el móvil»: muchas funciones de IA viven en el ordenador o la web; desde el móvil, con Sunlogin/ToDesk o la app del producto, puedes «pilotar» esas tareas de IA en trayectos o fuera de casa. Regla: cálculo pesado y documentos largos, al PC; ver el avance y retocar un par de frases, al móvil.
  \item Privacidad: antes de subir contratos, documentos de identidad o datos de clientes, revisa las condiciones de privacidad y la normativa de la empresa; con material corporativo, prioriza las herramientas que indique tu organización.
\end{itemize}

\booksubsection{Escritorio animado: fondos dinámicos de miHoYo (móvil + PC)}

\begin{itemize}
  \item Desktop Live Wallpaper (miHoYo) Público: fans de personajes de miHoYo que quieren fondos animados gratis. Puntos fuertes: app gratuita de fondos dinámicos de miHoYo, disponible en móvil y PC; personajes en Live2D/animaciones como fondo de escritorio y, en algunos modelos, también de pantalla de bloqueo.
  \item Batería y rendimiento: los fondos animados gastan más batería y GPU que una imagen fija; con poca batería o en equipos antiguos, úsalo solo al cargar o baja la frecuencia de imágenes si la app lo permite.
  \item Fuente de instalación: prioriza la web oficial o tiendas de apps fiables y evita versiones «crackeadas» con troyanos; en PC, revisa la ruta de instalación y si hace falta el inicio automático.
\end{itemize}

\booksubsection{Envío rápido en red local: LocalSend / Feige / AirDrop}

\begin{itemize}
  \item LocalSend Público: mezclas de Android/Windows/iOS/Mac/Linux que quieren enviar entre plataformas sin pasar por la nube. Puntos fuertes: código abierto y conexión directa en la misma red; fotos, vídeos y documentos por el mismo WiFi, sin cables ni servidores del fabricante; la velocidad depende del router.
  \item Feige (IPMsg) Público: oficinas con varios Windows y escenarios clásicos de LAN. Puntos fuertes: veterana herramienta de transferencia y chat en red local, ideal para archivos grandes en intranet; descárgala de fuentes fiables y respeta la política de TI de tu empresa.
  \item AirDrop Público: usuarios de iPhone, iPad y Mac. Puntos fuertes: compartir con un gesto en el ecosistema Apple; descubrimiento por Bluetooth y transferencia por WiFi directo, sin gastar datos; puedes limitarlo a «Solo contactos» para evitar molestias.
  \item Otras opciones: Xiaomi Share, Huawei Share, vivo Share, etc., van muy rápido entre dispositivos de la misma marca; el «Archivos de WeChat» es cómodo pero limita el tamaño y los dispositivos; para archivos grandes, mejor LocalSend o AirDrop.
  \item Consejos de uso: ambos en el mismo WiFi; en WiFi público, evita archivos sensibles y mejor un punto de acceso personal; al terminar, comprueba en el receptor que el archivo esté completo.
\end{itemize}

\booksubsection{Portapapeles entre dispositivos: KDE Connect y otras utilidades}

\begin{itemize}
  \item Portapapeles universal de Apple Público: usuarios de iPhone + Mac. Puntos fuertes: con la misma Apple ID y Handoff activado, copias en el móvil y pegas en el Mac (o al revés); texto y algunos archivos.
  \item KDE Connect Público: Android + Windows/Linux/Mac, con portapapeles de código abierto, espejo de notificaciones y entrada remota. Puntos fuertes: al emparejar en la misma red, sincroniza el portapapeles, envía notificaciones del móvil al PC y usa el móvil como mando o ratón; los datos no pasan por cuentas de terceros.
  \item Utilidades de marca «móvil-PC»: Huawei Multi-Screen Collaboration / Super Device, Xiaomi MiMo, vivo Share, OPPO Connect, Honor Connect, etc.; suelen traer portapapeles multi-pantalla y arrastrar archivos, con la mejor experiencia dentro de la misma marca.
  \item Canal oficial Windows + Android: «Enlace con Windows / Phone Link» (evolución de Your Phone) permite ver notificaciones y fotos del móvil en el PC y, en algunos casos, copiar y pegar; en Samsung y otras marcas está más integrado.
  \item Seguridad: el portapapeles puede retener contraseñas o códigos; en ordenadores compartidos, cuidado con la sincronización global y evita pegar de más en campos de contraseña; algunos teclados virtuales también leen el portapapeles: revisa permisos.
\end{itemize}

\booksubsection{Códigos dinámicos 2FA: Aegis / Microsoft Authenticator}

\begin{itemize}
  \item Por qué un autenticador y no SMS: el código por SMS puede secuestrarse con una SIM duplicada; el TOTP se genera localmente por algoritmo y hora, sin depender de la cobertura, y resiste mejor el phishing y la intercepción (aun así, cuidado con malware y capturas).
  \item Aegis Authenticator Público: usuarios de Android que prefieren código abierto y control local. Puntos fuertes: gratuito y abierto, con copias cifradas e importar/exportar; interfaz limpia y desbloqueo de tokens con huella o rostro.
  \item Microsoft Authenticator Público: cuentas Microsoft, Office 365 o corporativas; en iOS y Android. Puntos fuertes: aprobación por notificación, monitorización de contraseñas y copia de cuentas (según versión); muy integrado con Windows e identidad corporativa.
  \item Google Authenticator, 2FAS y otros: también habituales dentro de cada ecosistema; lo clave es que tengan copia de seguridad cifrada, para no perder los tokens al cambiar de móvil.
  \item Imprescindible al cambiar de móvil: migra desde el teléfono antiguo o usa los «códigos de respaldo/claves de recuperación» del servicio; imprime los códigos de respaldo en papel y guárdalos en dos sitios, no solo en el mismo móvil.
  \item Con el gestor de contraseñas: algunas bóvedas de contraseñas también guardan códigos TOTP, lo cual es cómodo pero concentra el riesgo; para cuentas de alto valor se recomienda repartir funciones entre «gestor de contraseñas + autenticador independiente». Más detalles en el Capítulo 14.
\end{itemize}

\booksubsection{Combinación recomendada: cómo encadenarlas en el día a día}

\textbf{Ejemplo de flujo de trabajo práctico}\par

\begin{itemize}
  \item Por la mañana: revisas la agenda del día en TickTick echas un vistazo al tiempo y a las alertas enciendes el ordenador con Trae o tus documentos y dejas el móvil a un lado.
  \item Durante el día: los documentos en papel los escaneas a PDF con una app de escaneo y los envías al ordenador con LocalSend o AirDrop; los textos largos los redactas en el ordenador con herramientas de ofimática con IA.
  \item Si sales: te conectas a tu ordenador con ToDesk o Sunlogin para seguir con la tarea; para iniciar sesión en cuentas importantes usas los códigos dinámicos de Aegis o Microsoft Authenticator.
  \item Por la noche: marcas como completadas las tareas en TickTick; si lo necesitas, cambias el fondo de pantalla del escritorio para relajarte un rato.
\end{itemize}

\begin{tcbwarn}{Reglas de oro para el control remoto y la seguridad de cuentas}
\begin{itemize}
  \item Si alguien te pide instalar un software de control remoto cuelga el teléfono y busca tú mismo el servicio de atención al cliente oficial.
  \item No le digas a nadie los códigos de verificación ni las contraseñas dinámicas, ni siquiera a quien dice ser «mensajería» o «atención al cliente».
  \item Antes de controlar el ordenador de la empresa a distancia o instalar software de IA, pregunta primero a IT si está permitido.
  \item Extrema la precaución al transferir archivos o iniciar sesión a distancia en redes públicas.
\end{itemize}

\end{tcbwarn}

\begin{tcbpractice}{Ejercicio práctico: completa una vez toda la cadena «entre dispositivos»}
\begin{enumerate}
  \item Conecta el móvil y el ordenador a la misma red Wi-Fi y envíate una foto entre ambos con LocalSend (o AirDrop).
  \item Crea en TickTick la tarea «Hoy a las 20:00: ordenar los documentos escaneados» y comprueba que se ve tanto en el móvil como en el ordenador.
  \item Activa la 2FA en tu correo o en WeChat (si lo permite), vincúlala con Aegis o Microsoft Authenticator y guarda los códigos de respaldo.
  \item Instala Sunlogin o ToDesk en tu propio ordenador, conéctate una vez con éxito desde tu móvil y luego cierra la sesión o refuerza la contraseña: familiarízate con el proceso sin dejar puertas traseras de seguridad.
\end{enumerate}

\end{tcbpractice}
